\section{The Vershik successor: the serial yarn order}

The third Bratteli--Vershik piece is the adic successor --- the serial yarn order, the dyadic
odometer, green in \texttt{Obsmetrologia.Closure.Odometer}. With it the Bratteli--Vershik system of a
self-similar textile is complete in-tree: edges (\cref{thm:port-growth}), levels-as-time
(\cref{thm:weaving-advances-time}), and the serial successor.

\begin{proposition}[The Vershik adic successor]\label{thm:vershik}\lean{Obsmetrologia.Odometer.vershik}\leanok\uses{thm:port-growth}
The next path in the serial order, the little-endian ``$+1$ with carry'' odometer: a not-yet-maximal
edge advances and stops; a maximal edge resets and carries up a level. The serial yarn order ---
finish the stitch, carry to the next row.
\end{proposition}

\begin{theorem}[The successor is the serial $+1$]\label{thm:vershik-succ}\lean{Obsmetrologia.Odometer.vershik_succ}\leanok\uses{thm:vershik}
On any path not yet maximal, the adic successor increments the serial index:
$\texttt{toNat}(\texttt{vershik}\,p) = \texttt{toNat}\,p + 1$. The serial yarn order is the integer
successor until a row maxes out.
\end{theorem}

\begin{theorem}[The maximal path wraps to the minimal]\label{thm:vershik-wrap}\lean{Obsmetrologia.Odometer.vershik_replicate_true}\leanok\uses{thm:vershik}
The all-maximal path of any height carries past the top and returns to the all-minimal path --- the
adic map's one discontinuity, where the serial order closes into a cycle.
\end{theorem}

\begin{theorem}[On two levels the successor is a 4-cycle]\label{thm:vershik-cycle}\lean{Obsmetrologia.Odometer.vershik_two_cycle}\leanok\uses{thm:vershik-succ, thm:vershik-wrap}
Iterating \texttt{vershik} four times returns any two-level path to itself: the minimal dyadic
odometer, every path reached --- the Bratteli--Vershik minimality of the two-row self-similar tower.
\end{theorem}

\begin{proposition}[The dyadic tower]\label{thm:dyadic-tower}\lean{Obsmetrologia.Odometer.dyadicTower}\leanok\uses{thm:vershik}
The dyadic tower as an inverse system: stage $n$ is the length-$n$ paths,
restriction is prefix truncation, functoriality by the take laws. Its
threads are the hull of the tower.
\end{proposition}

\begin{theorem}[The dyadic hull is Cantor space]\label{thm:dyadic-hull-cantor}\lean{Obsmetrologia.Odometer.dyadicHull_is_cantor}\leanok\uses{thm:dyadic-tower}
Reading a Cantor point's prefix thread recovers the point, and every thread
is the thread of its reads: the inverse limit $\varprojlim \{0,1\}^n$ is
$2^{\mathbb{N}}$, in coordinates.
\end{theorem}

\begin{theorem}[The odometer is a tower endomorphism]\label{thm:vershik-tower-endomorphism}\lean{Obsmetrologia.Odometer.vershikStage_proj}\leanok\uses{thm:dyadic-tower, thm:vershik-succ}
The Vershik successor commutes with prefix truncation
(\texttt{vershik\_take}: the carry propagates left to right), so the
successor acts stage-wise on the tower and descends to the Cantor hull
(\texttt{vershikHull}) --- the odometer as a dynamic on the inverse limit,
not only on finite paths.
\end{theorem}

\begin{theorem}[The stage is one full cycle]\label{thm:vershik-cycle-full}\lean{Obsmetrologia.Odometer.vershikIter_cycle}\leanok\uses{thm:vershik-succ, thm:vershik-wrap}
The $2^L$-fold successor is the identity on $L$-level paths, and no smaller
positive iterate fixes any path: in serial coordinates the $m$-fold iterate
adds $m$ modulo $2^L$ (\texttt{toNat\_vershikIter}, stated as an exact
equation with a wrap count), the serial index separates paths of equal level
(\texttt{eq\_of\_toNat\_eq}), and the index is bounded by $2^L$
(\texttt{toNat\_lt}). Generalizes \cref{thm:vershik-cycle} from two levels to
all levels.
\end{theorem}

\begin{theorem}[The periodic-point law of a stage]\label{thm:vershik-periodic-law}\lean{Obsmetrologia.Odometer.vershikIter_fixed_iff}\leanok\uses{thm:vershik-cycle-full}
The $m$-fold successor fixes an $L$-level path exactly when $2^L \mid m$.
Each stage is a single $2^L$-cycle, so every iterate fixes all $2^L$ paths
or none: the fixed-point count of the $m$-th iterate at stage $L$ is
$2^L \cdot [2^L \mid m]$ --- the complete periodic-orbit data of the finite
stages.
\end{theorem}

\begin{theorem}[The hull has no periodic points]\label{thm:hull-aperiodic}\lean{Obsmetrologia.Odometer.vershikHullIter_ne_self}\leanok\uses{thm:vershik-periodic-law, thm:vershik-tower-endomorphism}
No positive iterate of the odometer fixes any thread of the Cantor hull: a
fixed thread would be fixed at every stage, and any stage with $m < 2^n$
refuses by the periodic-point law. Every finite stage is periodic; the
inverse limit has empty periodic spectrum --- aperiodicity of the hull
dynamic as a theorem about the tower's own carrier, with no regime data
asserted.
\end{theorem}

\begin{theorem}[The periodicity dichotomy of the tower]\label{thm:tower-periodicity-dichotomy}\lean{Obsmetrologia.Odometer.tower_periodicity_dichotomy}\leanok\uses{thm:vershik-periodic-law, thm:hull-aperiodic}
The bundle: at every finite stage the fixed-point counts are $2^n$ at
multiples of $2^n$ and $0$ elsewhere; on the hull all counts vanish. The
periodic-orbit spectrum of stage $n$ is the single orbit length $2^n$ ---
the tower sees only the prime $2$ --- and in the limit the spectrum empties.
\end{theorem}

\begin{theorem}[Each stage is a single orbit]\label{thm:vershik-stage-transitive}\lean{Obsmetrologia.Odometer.vershikIter_reaches}\leanok\uses{thm:vershik-cycle-full}
Any two paths of the same level are connected by fewer than $2^L$ successor
steps: the stage orbit is one cycle through all $2^L$ paths, so the stage
orbit graph has $2^L$ vertices, $2^L$ edges, and one component.
\end{theorem}

\begin{theorem}[The hull odometer is minimal and aperiodic]\label{thm:vershik-hull-minimal}\lean{Obsmetrologia.Odometer.vershikHull_minimal_aperiodic}\leanok\uses{thm:vershik-stage-transitive, thm:hull-aperiodic, thm:vershik-tower-endomorphism}
Every orbit of the hull odometer is dense --- for any two threads and any
stage, some iterate carries the first to agree with the second at that stage
(every cylinder is visited) --- and no positive iterate has a fixed point.
The module header's ``canonical minimal Cantor system'' is now a theorem,
not prose: minimality by \texttt{vershikHull\_orbit\_dense}, aperiodicity by
\cref{thm:hull-aperiodic}.
\end{theorem}

\begin{theorem}[The serial chart and the `Fin` avatar]\label{thm:vershik-serial-chart}\lean{Obsmetrologia.Odometer.vershikFinIter_cycle}\leanok\uses{thm:vershik-cycle-full, thm:vershik-stage-transitive}
The stage orbit is its own coordinate chart: \texttt{ofNat n v} (the $v$-th
successor step from the minimal path) is a two-sided inverse to the serial
index \texttt{toNat} (\texttt{toNat\_ofNat}, \texttt{ofNat\_toNat}), with no
modular arithmetic stated anywhere --- the chart is the dynamics. Conjugating
gives the successor's $\mathrm{Fin}(2^n)$ avatar \texttt{vershikFin}, whose
$2^n$-fold iterate is the identity (\texttt{vershikFinIter\_cycle}), by pure
transport of \cref{thm:vershik-cycle-full}.
\end{theorem}

\begin{theorem}[The winding exchange]\label{thm:vershik-winding-exchange}\notready\uses{thm:tower-periodicity-dichotomy, thm:vershik-stage-transitive, thm:coordination-h1rank}
The wrap count $k$ of \texttt{toNat\_vershikIter} is unique
(\texttt{toNat\_vershikIter\_unique}), so an $m$-step run at stage $n$ has a
well-defined winding number, and a periodic return spends exactly $m = k
\cdot 2^n$ steps for $k$ windings. Identify this with the coordination
cohomology: derive the stage orbit graph's cover-nerve counts ($2^n$
vertices and edges from \cref{thm:vershik-stage-transitive}, one component
from transitivity, needing the missing component counter), giving
$\check H^1 = 1$ --- each stage carries exactly one independent cycle, one
Landauer bit, and the winding number is the pairing of the run against that
one $\check H^1$ generator. The open work: a derived component counter for
orbit graphs, the cover-nerve instance, and the pairing statement.
\end{theorem}

\begin{theorem}[The Artin--Mazur zeta of the tower]\label{thm:tower-artin-mazur-zeta}\notready\uses{thm:tower-periodicity-dichotomy}
Assemble the fixed-point counts of \cref{thm:tower-periodicity-dichotomy}
into the Artin--Mazur zeta $\zeta_n(t) = \exp\bigl(\sum_m
|\mathrm{Fix}|\, t^m/m\bigr)$ over the library's own rational/formal-series
stack: stage $n$ has $\zeta_n(t) = (1 - t^{2^n})^{-1}$, rational with all
inverse zeros of modulus $1$, and the hull has $\zeta = 1$. This is the
trace-formula exchange in miniature: local periodic-orbit counts against a
spectral side, the exchange whose archimedean case is the Riemann explicit
formula and whose finite-field case (Frobenius eigenvalues on $H^1$, weights
on the critical line) is the proved Riemann hypothesis of Weil--Deligne. The
open work: a formal-power-series carrier over the Bishop rationals, the
rationality statement, and the identification of the stage spectra as the
$2^n$-th roots of unity.
\end{theorem}

\begin{proposition}[The Bernoulli cylinder law]\label{thm:bernoulli-cylinder-law}\lean{Obsmetrologia.Odometer.bernoulliLaw}\leanok\uses{thm:dyadic-tower}
A cylinder law assigns each cylinder word a pre-rational mass, with a
cylinder's mass cross-multiplicatively equal to the sum of its two
children's and unit mass on the empty word; the laws are stated at the
\texttt{PreRat} setoid level. The uniform assignment $1/2^n$ per depth-$n$
cylinder is such a law (\texttt{bernoulliMass\_split}, unit total), and
each stage distribution is uniform (\texttt{bernoulliMass\_uniform}). The
quotient-level face is \texttt{bernoulli\_split} /
\texttt{bernoulli\_total} over the derived rationals.
\end{proposition}

\begin{theorem}[The odometer preserves the Bernoulli law]\label{thm:bernoulli-vershik-invariant}\lean{Obsmetrologia.Odometer.bernoulliLaw_vershik_invariant}\leanok\uses{thm:bernoulli-cylinder-law, thm:vershik}
The adic successor carries each cylinder to a cylinder of equal Bernoulli
mass --- an equality of pre-rationals, by length preservation
(\texttt{bernoulliMass\_vershik} via \texttt{vershik\_length}) --- so the
Bernoulli law is invariant under the odometer.
\end{theorem}

\begin{theorem}[Stage unique ergodicity]\label{thm:odometer-stage-unique-ergodicity}\notready\uses{thm:bernoulli-vershik-invariant, thm:vershik-stage-transitive}
Each stage is a single $2^n$-cycle (\cref{thm:vershik-stage-transitive}),
so a vershik-invariant stage distribution is constant along the orbit: the
uniform stage mass of \cref{thm:bernoulli-cylinder-law} is the only
invariant one. The open work: the constancy argument over the derived
rationals, and the inverse-limit statement identifying the Bernoulli law
as the hull odometer's unique invariant law.
\end{theorem}
